\documentclass{article}

\IfFileExists{styles/main.sty}{%
  \usepackage{styles/main}%
}{%
  \usepackage{../../styles/main}%
}

\usepackage{xurl}
\usepackage{cleveref}

\IfFileExists{references.bib}{%
  \addbibresource{references.bib}%
}{%
  \addbibresource{pws/2/references.bib}%
}

\IfFileExists{figures/}{%
  \graphicspath{{./figures/}}%
}{%
  \graphicspath{{./pws/2/figures/}}%
}

\begin{document}

% ==================== TITLE PAGE ====================
\IfFileExists{docs/title.pdf}{%
  \includepdf[pages=-]{docs/title.pdf}%
}{%
  \includepdf[pages=-]{pws/2/docs/title.pdf}%
}

\selectlanguage{ukrainian}

\crefname{figure}{рис.}{рис.}
\Crefname{figure}{Рисунок}{Рисунки}
\crefname{table}{табл.}{табл.}
\Crefname{table}{Таблиця}{Таблиці}
\crefname{section}{розд.}{розд.}
\Crefname{section}{Розділ}{Розділи}
\crefname{equation}{формула}{формули}
\Crefname{equation}{Формула}{Формули}
\renewcommand{\crefrangeconjunction}{~--~}
\renewcommand{\crefpairconjunction}{ та~}
\renewcommand{\crefmiddleconjunction}{, }
\renewcommand{\creflastconjunction}{ та~}

\titleformat{\section}[block]
{\filcenter\bfseries\fontsize{16pt}{24pt}\selectfont}
{}
{0pt}
{\MakeUppercase}
\titlespacing*{\section}{0pt}{20pt}{6pt}

% ==================== CONTENT ====================
\section{ПРАКТИЧНА РОБОТА № 2}

\textbf{Тема:} Оформлення списку літератури з використанням систем управління бібліографічною інформацією.

\textbf{Мета:} Встановити, налаштувати та навчитися використовувати безкоштовний бібліографічний менеджер Zotero для оформлення списку літератури у наукових роботах.

\subsection{Постановка завдання}
Інсталювати бібліографічний менеджер Zotero, навчитися ефективно його використовувати.

\subsection{Порядок виконання роботи}

\subsubsection{Текст до теми БКР «Розробка системи керування рухом нижніх кінцівок двоногого робота»}

Керування рухом нижніх кінцівок двоногого робота ускладнюється нелінійною динамікою, розривними контактами з поверхнею та обмеженнями на крутні моменти приводів. Для розв'язання задачі динамічної рівноваги застосовують модельно-прогнозувальне керування (MPC) на основі критерію точки нульового моменту (ZMP)~\cite{xuZMPAwareTaskFormulation2026}, що оптимізує контактні сили на скінченному горизонті. Перенесення отриманих сил на суглоби здійснює контролер усього тіла (WBC) із контролем допустимості контактів~\cite{zhangFeasibilityawareWholebodyControl2026}, який запобігає ковзанню стоп і стабілізує позу під час руху складним рельєфом.

Для підвищення робастності до зовнішніх збурень модельно-орієнтовані алгоритми доповнюють глибоким навчанням з підкріпленням (DRL). Автоматичне формування функцій винагороди суттєво спрощує перенесення політик із симуляції на фізичного робота (sim-to-real)~\cite{luDeepReinforcementLearning2026}. Своєю чергою, злиття мультимодальних сенсорних потоків (зору та пропріоцепції)~\cite{jiLearningbasedLocomotionControl2025} і використання перцептивних внутрішніх моделей~\cite{longLearningHumanoidLocomotion2025} дозволяють адаптувати параметри кроку до невідомих перешкод та геометрії опорної поверхні в реальному часі.

Критичним аспектом практичної реалізації системи керування є надання формальних гарантій безпеки для уникнення падінь робота. Використання ризик-чутливих безакторних стратегій на основі опуклої оптимізації~\cite{zhangRisksensitiveActorfreePolicy2023} дає змогу обмежити ймовірність виходу за межі допустимих станів. Комплексна інтеграція MPC, WBC та безпечного машинного навчання формує цілісну й енергоефективну архітектуру керування нижніми кінцівками двоногого робота.

\subsubsection{Налаштування Zotero}
На офіційному веб-порталі бібліографічного менеджера Zotero~\cite{zotero} було зареєстровано персональний обліковий запис з іменем користувача \textit{AndriiOsidach} (\cref{fig:web_zotero}). Це дозволило отримати доступ до хмарного сховища Zotero Cloud та онлайн-бібліотеки (\textit{Web Library}), що забезпечує збереження бібліографічної бази й доступ до неї з будь-якого пристрою через браузер. У структурі особистої бібліотеки («My Library») було створено тематичну колекцію \texttt{Science} для впорядкування наукових публікацій до цієї дисципліни.

Для локальної роботи на комп'ютері під управлінням операційної системи Linux було завантажено та встановлено десктопну версію програми Zotero Standalone (\cref{fig:local_zotero}).

У налаштуваннях десктопного додатку («Preferences $\rightarrow$ Sync») було виконано авторизацію з використанням створеного облікового запису \textit{AndriiOsidach} та увімкнено автоматичну двосторонню синхронізацію даних із хмарою Zotero. У результаті структуру персональної бібліотеки та колекцію \texttt{Science} було успішно синхронізовано між локальною програмою та веб-бібліотекою (\cref{fig:web_zotero,fig:local_zotero}).

\begin{figure}[h]
  \centering
  \includegraphics[width=0.88\linewidth]{web_zotero.jpg}
  \caption{Веб-інтерфейс бібліотеки користувача у бібліографічному менеджері Zotero.}
  \label{fig:web_zotero}
\end{figure}

\begin{figure}[h]
  \centering
  \includegraphics[width=0.88\linewidth]{local_zotero.jpg}
  \caption{Робоче вікно десктопної версії програми Zotero у середовищі Linux.}
  \label{fig:local_zotero}
\end{figure}

\clearpage
\subsubsection{Джерела у Zotero}

За допомогою браузерного розширення Zotero Connector для Microsoft Edge безпосередньо з веб-сторінок видавництв до тематичної колекції \texttt{Science} було додано наукові статті з MDPI (журнал \textit{Symmetry})~\cite{xuZMPAwareTaskFormulation2026}, Springer (журнал \textit{Autonomous Robots})~\cite{zhangFeasibilityawareWholebodyControl2026} та Elsevier (журнал \textit{Biomimetic Intelligence and Robotics})~\cite{jiLearningbasedLocomotionControl2025}. Під час імпорту розширення автоматично розпізнало цифровий ідентифікатор DOI на сторінках публікацій і на його основі підтягнуло всю необхідну бібліографічну інформацію, метадані та повнотекстові копії статей у форматі PDF.

Джерело з відкритого порталу CEUR-WS.org авторства R. Zhang та J. Sjölund~\cite{zhangRisksensitiveActorfreePolicy2023} було додано до бібліотеки вручну, оскільки праця опублікована у збірнику CEUR Workshop Proceedings (том 3505) та не містить цифрового ідентифікатора DOI. У середовищі десктопної програми Zotero було створено новий бібліографічний запис типу «Conference Paper» та власноруч заповнено всі необхідні інформаційні поля (\cref{fig:manual_ceus_ws}): назву статті, авторів, редакторів збірника, повну назву конференції, видавця CEUR, дату оприлюднення, номер тому, ISSN, а також гіперпосилання на електронну версію матеріалів (URL).

\begin{figure}[h]
  \centering
  \includegraphics[width=0.88\linewidth]{manual_ceus-ws.png}
  \caption{Ручне додавання бібліографічного джерела без DOI (CEUR-WS.org) у десктопній версії програми Zotero.}
  \label{fig:manual_ceus_ws}
\end{figure}

Два інші джерела з міжнародних баз IEEE та PubMed були додані за допомогою вбудованої функції десктопної версії Zotero «Add Item(s) by Identifier» шляхом введення їхніх цифрових ідентифікаторів DOI. У такий спосіб було імпортовано матеріали виступу на конференції IEEE ICRA 2025~\cite{longLearningHumanoidLocomotion2025} та статтю у науковому журналі \textit{Research}, проіндексовану у базі даних PubMed~\cite{luDeepReinforcementLearning2026}. Бібліографічний менеджер автоматично виконав онлайн-запит до глобальних реєстрів ідентифікаторів, імпортував структуровані метадані публікацій та завантажив пов'язані електронні копії у персональне сховище.

% ==================== BIBLIOGRAPHY ====================
\clearpage
\nocite{*}
\printbibliography[heading=bibintoc,title={СПИСОК ВИКОРИСТАНИХ ДЖЕРЕЛ}]

\end{document}
